%%%%%%%%%%%%%%%%%%%%%%%%%%%%%%%%%%%%%%%%%%%%%%%%%%%%%%%%%%%%%%%%%%%%%%%%
%                                                                      %
%     File: Thesis_Nomenclature.tex                                    %
%     Tex Master: Thesis.tex                                           %
%                                                                      %
%     Author: Andre C. Marta                                           %
%     Last modified : 29 Jun 2022                                      %
%                                                                      %
%%%%%%%%%%%%%%%%%%%%%%%%%%%%%%%%%%%%%%%%%%%%%%%%%%%%%%%%%%%%%%%%%%%%%%%%
%
% The definitions can be placed anywhere in the document body
% and their order is sorted by <symbol> automatically when
% calling makeindex in the makefile
%
% The \glossary command has the following syntax:
%
% \glossary{entry}
%
% The \nomenclature command has the following syntax:
%
% \nomenclature[<prefix>]{<symbol>}{<description>}
%
% where <prefix> is used for fine tuning the sort order,
% <symbol> is the symbol to be described, and <description> is
% the actual description.

% ----------------------------------------------------------------------
% Roman symbols [r]
\nomenclature[R]{$\mathbf{q}$}{Mobile robot pose.}
\nomenclature[R]{$\mathbf{p}$}{Mobile robot position.}
\nomenclature[R]{$\mathbf{\chi}$}{Mobile robot orientation.}
\nomenclature[R]{$\mathbf{v}$}{Mobile robot velocity.}
\nomenclature[R]{$\mathcal{B}, \mathcal{W}$}{Body and world frames, respectively.}
\nomenclature[R]{$\mathbf{R}$}{Rotation matrix.}

% ----------------------------------------------------------------------
% Greek symbols [g]
\nomenclature[G]{$\phi, \; \theta, \; \psi$}{Roll, pitch, and yaw angles, respectively.}

% ----------------------------------------------------------------------
% Subscripts [s]
%\nomenclature[S]{$x,y,z$}{Cartesian components}

% ----------------------------------------------------------------------
% Supercripts [t]
%\nomenclature[T]{T}{Transpose}

